\documentclass[11pt]{article}

\usepackage[utf8]{inputenc}
\usepackage[T1]{fontenc}
\usepackage{lmodern}
\usepackage{geometry}
\usepackage{amsmath,amssymb,amsfonts,amsthm,mathtools}
\usepackage{bm}
\usepackage{enumitem}
\usepackage{microtype}
\usepackage{hyperref}
\usepackage{titlesec}
\usepackage{booktabs}
\usepackage{array}
\usepackage{setspace}
\usepackage{mathrsfs}
\usepackage{physics}

\geometry{margin=1in}
\hypersetup{colorlinks=true, linkcolor=black, urlcolor=blue, citecolor=black}
\setlist[enumerate]{topsep=0.4em,itemsep=0.35em}
\setlist[itemize]{topsep=0.3em,itemsep=0.25em}
\titleformat{\section}{\large\bfseries}{\thesection.}{0.5em}{}
\titleformat{\subsection}{\normalsize\bfseries}{\thesubsection.}{0.5em}{}
\setstretch{1.06}

\newcommand{\Aether}{\AE{}ther}
\newcommand{\Aflow}{\AE{}ther-flow}
\newcommand{\TheoryName}{\Aflow{} Interpretation of Relativity}
\newcommand{\TheoryProgram}{\Aether{} / \Aflow{} framework}
\newcommand{\ObserverMetric}{g^{\mathrm{br}}}
\newcommand{\CandidateMetric}{g^{\mathrm{cand}}}
\newcommand{\CompletedMetric}{g^{\sharp}}
\newcommand{\RedefMetric}{\widehat g}
\newcommand{\BaseShift}{\Sigma^{\mathrm{IER}}}
\newcommand{\ResponseShift}{\Sigma^{\mathrm{resp}}}
\newcommand{\CompShift}{\Sigma^{\mathrm{cmp}}}
\newcommand{\BridgeShift}{\Sigma^{\mathrm{brg}}}
\newcommand{\ResidualCore}{\mathcal V_{\mu\nu}^{\mathrm{res}}}
\newcommand{\ResidualDec}{\mathcal D_{\mu\nu}^{\mathrm{res}}}
\newcommand{\MixQuad}{\mathcal M_{\mu\nu}^{\mathrm{mix}}}
\newcommand{\RespQuad}{\mathcal Q_{\mu\nu}^{\mathrm{resp}}}
\newcommand{\DeltaTCand}{\Delta T_{\mu\nu}^{\mathrm{cand}}}
\newcommand{\ReducedEin}{\mathcal L_{\ObserverMetric}}
\newcommand{\GaugeBook}{\mathcal G_{\ObserverMetric}}
\newcommand{\CompMix}{\mathcal M_{\mu\nu}^{\mathrm{cmp}}}
\newcommand{\CompQuad}{\mathcal Q_{\mu\nu}^{\mathrm{cmp}}}
\newcommand{\DeltaTComp}{\Delta T_{\mu\nu}^{\mathrm{cmp}}}
\newcommand{\ObsBurden}{\mathfrak B_{\mu\nu}^{\mathrm{obs}}}
\newcommand{\ObserverClass}[1]{\left[#1\right]_{\mathrm{obs}}}
\newcommand{\GaugeExact}{\mathfrak G_{\mu\nu}^{\mathrm{ex}}}
\newcommand{\ConstraintExact}{\mathfrak C_{\mu\nu}^{\mathrm{ex}}}
\newcommand{\ProtoFunctional}{\mathscr F_{\mathrm{proto}}}
\newcommand{\StemFunctional}{\mathscr F_{\mathrm{sel}}}
\newcommand{\SharpFunctional}{\mathscr F_{\mathrm{sharp}}}
\newcommand{\RootTensorClass}{\mathscr U_{\mathrm{car}}}
\newcommand{\RootRelabelClass}{\mathscr U_{\mathrm{cur}}}
\newcommand{\RootFreeClass}{\mathscr U_{\mathrm{hom}}}
\newcommand{\RootTensor}{\mathfrak U}
\newcommand{\RootRelabel}{\mathfrak V}
\newcommand{\RootFree}{\mathfrak Q}
\newcommand{\TensorSymProj}{\Pi_{\mathrm{sym,car}}}
\newcommand{\CurSymProj}{\Pi_{\mathrm{sym,cur}}}
\newcommand{\HomSymProj}{\Pi_{\mathrm{sym,hom}}}
\newcommand{\TensorSymGroup}{G_{\mathrm{car}}}
\newcommand{\CurSymGroup}{G_{\mathrm{cur}}}
\newcommand{\HomSymGroup}{G_{\mathrm{hom}}}
\newcommand{\TensorSymAct}{\mathscr U_{\mathrm{car}}}
\newcommand{\CurSymAct}{\mathscr U_{\mathrm{cur}}}
\newcommand{\HomSymAct}{\mathscr U_{\mathrm{hom}}}
\newcommand{\TensorCompat}{\mathcal C_{\mathrm{car}}}
\newcommand{\CurCompat}{\mathcal C_{\mathrm{cur}}}
\newcommand{\HomCompat}{\mathcal C_{\mathrm{hom}}}
\newcommand{\TensorCmpProj}{\Pi_{\mathrm{cmp,car}}}
\newcommand{\CurCmpProj}{\Pi_{\mathrm{cmp,cur}}}
\newcommand{\HomCmpProj}{\Pi_{\mathrm{cmp,hom}}}
\newcommand{\TensorCurrent}{\mathfrak c_{\mathrm{car}}}
\newcommand{\CurCurrent}{\mathfrak c_{\mathrm{cur}}}
\newcommand{\HomCurrent}{\mathfrak c_{\mathrm{hom}}}
\newcommand{\TensorResp}{\mathcal R_{\mathrm{sel,car}}}
\newcommand{\CurResp}{\mathcal R_{\mathrm{sel,cur}}}
\newcommand{\HomResp}{\mathcal R_{\mathrm{sel,hom}}}
\newcommand{\TensorEqClass}{\mathscr E_{\mathrm{car}}}
\newcommand{\CurEqClass}{\mathscr E_{\mathrm{cur}}}
\newcommand{\HomEqClass}{\mathscr E_{\mathrm{hom}}}
\newcommand{\TensorEqRead}{\mathcal A_{\mathrm{car}}}
\newcommand{\CurEqRead}{\mathcal A_{\mathrm{cur}}}
\newcommand{\HomEqRead}{\mathcal A_{\mathrm{hom}}}
\newcommand{\TensorAmbientHess}{\mathbb H_{\mathrm{car}}}
\newcommand{\CurAmbientHess}{\mathbb H_{\mathrm{cur}}}
\newcommand{\HomAmbientHess}{\mathbb H_{\mathrm{hom}}}
\newcommand{\TensorEqKerProj}{\widetilde\Pi_{\mathrm{mic,car}}}
\newcommand{\CurEqKerProj}{\widetilde\Pi_{\mathrm{mic,cur}}}
\newcommand{\HomEqKerProj}{\widetilde\Pi_{\mathrm{mic,hom}}}
\newcommand{\TensorAmbientEnergy}{\mathcal U_{\mathrm{car}}}
\newcommand{\CurAmbientEnergy}{\mathcal U_{\mathrm{cur}}}
\newcommand{\HomAmbientEnergy}{\mathcal U_{\mathrm{hom}}}
\newcommand{\TensorEqManifold}{\mathcal N_{\mathrm{car}}}
\newcommand{\CurEqManifold}{\mathcal N_{\mathrm{cur}}}
\newcommand{\HomEqManifold}{\mathcal N_{\mathrm{hom}}}
\newcommand{\TensorTheta}{\Theta_{\mathrm{car}}}
\newcommand{\CurTheta}{\Theta_{\mathrm{cur}}}
\newcommand{\HomTheta}{\Theta_{\mathrm{hom}}}
\newcommand{\TensorCurrentBridge}{\mathcal B_{\mathrm{car}}^{\mathrm{cur}}}
\newcommand{\CurCurrentBridge}{\mathcal B_{\mathrm{cur}}^{\mathrm{cur}}}
\newcommand{\HomCurrentBridge}{\mathcal B_{\mathrm{hom}}^{\mathrm{cur}}}
\newcommand{\TensorEqBridge}{\mathcal B_{\mathrm{car}}^{\mathrm{eq}}}
\newcommand{\CurEqBridge}{\mathcal B_{\mathrm{cur}}^{\mathrm{eq}}}
\newcommand{\HomEqBridge}{\mathcal B_{\mathrm{hom}}^{\mathrm{eq}}}
\newcommand{\TensorShiftLift}{\mathcal S_{\mathrm{car}}^{\mathrm{cur}}}
\newcommand{\CurShiftLift}{\mathcal S_{\mathrm{cur}}^{\mathrm{cur}}}
\newcommand{\HomShiftLift}{\mathcal S_{\mathrm{hom}}^{\mathrm{cur}}}
\newcommand{\TensorEqLift}{\mathcal S_{\mathrm{car}}^{\mathrm{eq}}}
\newcommand{\CurEqLift}{\mathcal S_{\mathrm{cur}}^{\mathrm{eq}}}
\newcommand{\HomEqLift}{\mathcal S_{\mathrm{hom}}^{\mathrm{eq}}}
\newcommand{\ResidualRead}{\mathscr R_{\mathrm{res}}}
\newcommand{\CompletionRead}{\mathscr R_{\mathrm{cmp}}}
\newcommand{\MatterRead}{\mathscr R_{\mathrm{mat}}}

\newtheorem{definition}{Definition}
\newtheorem{lemma}{Lemma}
\newtheorem{proposition}{Proposition}
\newtheorem{remark}{Remark}
\newtheorem{corollary}{Corollary}
\newtheorem{theorem}{Theorem}

\title{\textbf{The \TheoryName{}}\\[0.4em]
\large Substrate-Response / Self-Response Charge-Polarization Candidate\\
\large Exact-Preservation Observer-Equation Continuity-Hessian\\
\large Bridge Theorem-Building Note}
\author{}
\date{}

\begin{document}

\maketitle

\begin{abstract}
The fixed-package strict-closure maximality theorem now determines the next honest move on the active one-metric charge-polarization line. The current observer package is exhausted: any strict observer-side closure upgrade on that fixed package must furnish either the stronger direct absorbability identity
\[
\ObserverClass{\ObsBurden}=0
\]
or an explicit same-package field-redefinition theorem. Therefore the next primary step is not another manuscript that merely reproduces the same downstream package.

The present note starts from the only recorded layer that is visibly connected to that missing theorem burden: the continuity-defect / ambient-Hessian interface, where the compatibility-current maps arise as exact continuity defects and the primitive response operators arise as pulled-back ambient Hessians. It then builds the smallest sharpened package at that interface. The symmetry actions, the visible \(\ProtoFunctional\) slot, and the one-metric branch data are kept fixed, but the linearized compatibility penalties are replaced by exact current-defect data and the quadratic response penalties are replaced by exact ambient-equilibrium data. The first theorem of the note proves that the recorded symmetry-charge / compatibility-response functional is exactly the quadratic truncation of this sharpened package. The second positive theorem proves that the preserved exact-preserving branch is already an exact sharpened branch, so branch locking on the active completed observer branch is now in hand.

That nonlinear sharpening alone is still not enough until the observer burden is actually tied back to the exact defect slots. The note therefore constructs observer-side sector readouts of \(\ResidualDec\), \(\CompMix+\CompQuad\), and \(-8\pi G_N\,\DeltaTComp\) directly from the ordered six-defect package, sums them to obtain the direct absorbability bridge identity \eqref{eq:direct_bridge}, and then applies branch locking to conclude
\[
\ObserverClass{\ObsBurden}=0
\]
on the active completed branch.

The claim boundary remains disciplined. This is a same-package observer-side closure theorem on the branch-locked continuity-Hessian sharpening, not a completed first-principles derivation of GR from substrate microphysics. The same-package rewritten-metric route is retained only as an optional alternative re-expression, not as the remaining front-line closure burden.
\end{abstract}

\section{Introduction}

The active manuscript sequence of \TheoryName{} has now crossed a clean internal threshold. On the fixed one-metric exact-preservation package, the completed observer equation has already been reduced to
\[
G_{\mu\nu}[\CompletedMetric]
=
8\pi G_N\,T_{\mu\nu}^{\mathrm{matter}}[\CompletedMetric,\psi]
+\GaugeBook(\CompShift)_{\mu\nu}
+\ObsBurden,
\]
with
\[
\ObsBurden
=
\ResidualDec+\CompMix+\CompQuad-8\pi G_N\,\DeltaTComp.
\]
The subsequent observer-equation notes then showed, in increasing strength, that the displayed sectors are benchmark-decoupled, exactly non-independent, not directly absorbable on the fixed package, and not yet removable by any presently recorded same-package field redefinition. The fixed-package maximality theorem made the project consequence explicit: at current scope there is no third same-package route.

That changes what counts as honest progress. A deeper-origin manuscript is no longer primary merely because it goes one layer lower while reproducing the same completed observer equation, the same completion solve, and the same one operative metric. To matter at the front line now, a sharpened package must change the observer-side data enough to prove one of the two missing theorems.

The only recorded interface that looks mathematically close to that burden is the continuity-defect / ambient-Hessian interface. At that interface, the compatibility-current maps are not primitive: they arise as continuity defects of deeper density readouts. Likewise the primitive response operators are not primitive: they are pullbacks of ambient equilibrium Hessians. Those are exactly the two places where the fixed package still remembers nonlinear data that the recorded observer line has not yet used.

The purpose of this note is therefore narrow and theorem-building in character. It does not draft another primary manuscript on the same fixed package. Instead it does four smaller things.
\begin{enumerate}
    \item It defines the smallest nonlinear sharpening of the continuity-Hessian interface that changes the observer-side data while keeping the same visible branch inputs.
    \item It proves that the already-recorded symmetry-charge / compatibility-response functional is exactly the quadratic truncation of that sharpening.
    \item It proves that the preserved exact-preserving branch is already an exact sharpened branch, so branch locking on the active completed observer branch holds on the sharpened package.
    \item It proves the direct-route factorization theorem on that branch-locked sharpening and leaves the same-package rewritten-metric theorem as a secondary alternative re-expression.
\end{enumerate}

\paragraph{Framework and claim boundary.}
\TheoryProgram{} denotes the broader \Aether{} / \Aflow{} framework built on the claims that the \Aether{} is the underlying four-dimensional substrate of reality, the \Aflow{} is its intrinsic ordered motion, observed three-dimensional space is the local experiential slice of that deeper substrate, S-time is the experienced order of change arising from matter, light, and the \Aflow{}, and observed expansion is the three-dimensional appearance of deeper four-dimensional ordered motion. Within that framework, \TheoryName{} denotes the exact-closure theory in which the effective gravitational dynamics are adopted to be exactly Einsteinian with universal matter coupling. In the active sequence, \emph{adoption} means use of that established relativistic dynamics without claiming substrate derivation, while \emph{derivation} is reserved for a first-principles recovery from explicit substrate variables.

The present note stays strictly inside that boundary. It does not claim a completed first-principles derivation of GR from substrate microphysics. Its narrower positive role is to prove the direct same-package bridge theorem on the smallest sharpened package that can honestly change the observer-side conclusion.

\section{Fixed Observer Bottleneck and the Recorded Interface}

\begin{definition}[Fixed one-metric observer package]
\label{def:fixed}
The present note keeps the following observer-side data fixed.
\begin{enumerate}
    \item The candidate and completed operative metrics are
    \begin{equation}
    \CandidateMetric
    =
    \ObserverMetric+\BaseShift+\ResponseShift,
    \qquad
    \CompletedMetric
    =
    \ObserverMetric+\BaseShift+\ResponseShift+\CompShift.
    \label{eq:metrics}
    \end{equation}
    \item The fixed completion solve is
    \begin{equation}
    \ReducedEin(\CompShift)=-\ResidualCore.
    \label{eq:solve}
    \end{equation}
    \item The completed observer equation is
    \begin{equation}
    G_{\mu\nu}[\CompletedMetric]
    =
    8\pi G_N\,T_{\mu\nu}^{\mathrm{matter}}[\CompletedMetric,\psi]
    +\GaugeBook(\CompShift)_{\mu\nu}
    +\ObsBurden,
    \label{eq:completed}
    \end{equation}
    with
    \begin{equation}
    \ObsBurden
    \equiv
    \ResidualDec+\CompMix+\CompQuad-8\pi G_N\,\DeltaTComp.
    \label{eq:burden}
    \end{equation}
    \item The already-recorded observer quotient mods out only the already-extracted null classes:
    \begin{enumerate}
        \item gauge-bookkeeping / gauge-exact tensors;
        \item constraint-exact elimination tensors.
    \end{enumerate}
    The observer-quotient class of a tensor \(X_{\mu\nu}\) is denoted \(\ObserverClass{X_{\mu\nu}}\).
    \item The fixed-package maximality theorem already recorded in the active sequence states that any admissible strict observer-side closure upgrade on this fixed package must furnish at least one of exactly two ingredients:
    \begin{enumerate}
        \item the direct absorbability identity \(\ObserverClass{\ObsBurden}=0\);
        \item an explicit same-package field-redefinition theorem.
    \end{enumerate}
\end{enumerate}
\end{definition}

\begin{definition}[Recorded continuity-Hessian interface]
\label{def:interface}
On the sharper symmetry-fixed sectors
\[
\operatorname{ran}\TensorSymProj\subset\RootTensorClass,\qquad
\operatorname{ran}\CurSymProj\subset\RootRelabelClass,\qquad
\operatorname{ran}\HomSymProj\subset\RootFreeClass,
\]
the recorded continuity-Hessian interface provides:
\begin{enumerate}
    \item Fr\'echet-smooth exact-preserving current maps
    \[
    \TensorCurrent,\qquad
    \CurCurrent,\qquad
    \HomCurrent
    \]
    with
    \begin{equation}
    \TensorCurrent(0)=0,\qquad
    \CurCurrent(0)=0,\qquad
    \HomCurrent(0)=0,
    \label{eq:current_zero}
    \end{equation}
    and linearizations
    \begin{equation}
    D\TensorCurrent(0)=\TensorCompat,\qquad
    D\CurCurrent(0)=\CurCompat,\qquad
    D\HomCurrent(0)=\HomCompat;
    \label{eq:current_linearization}
    \end{equation}
    \item exact-preserving ambient readouts
    \[
    \TensorEqRead:\operatorname{ran}\TensorSymProj\to\TensorEqClass,\qquad
    \CurEqRead:\operatorname{ran}\CurSymProj\to\CurEqClass,\qquad
    \HomEqRead:\operatorname{ran}\HomSymProj\to\HomEqClass;
    \]
    \item \(C^2\) ambient equilibrium energies
    \[
    \TensorAmbientEnergy,\qquad
    \CurAmbientEnergy,\qquad
    \HomAmbientEnergy
    \]
    with equilibrium manifolds
    \[
    \TensorEqManifold,\qquad
    \CurEqManifold,\qquad
    \HomEqManifold
    \]
    passing through the origin and satisfying
    \begin{equation}
    D\TensorAmbientEnergy(0)=0,\qquad
    D\CurAmbientEnergy(0)=0,\qquad
    D\HomAmbientEnergy(0)=0,
    \label{eq:energy_zero}
    \end{equation}
    together with second-variation operators
    \begin{equation}
    D^2\TensorAmbientEnergy(0)=\TensorAmbientHess,\qquad
    D^2\CurAmbientEnergy(0)=\CurAmbientHess,\qquad
    D^2\HomAmbientEnergy(0)=\HomAmbientHess;
    \label{eq:energy_hess}
    \end{equation}
    \item ambient equilibrium kernel projectors
    \[
    \TensorEqKerProj,\qquad
    \CurEqKerProj,\qquad
    \HomEqKerProj
    \]
    onto the tangent equilibrium directions at the origin, so the recorded primitive response operators satisfy
    \begin{equation}
    \TensorResp=\TensorEqRead^*\TensorAmbientHess\TensorEqRead,\qquad
    \CurResp=\CurEqRead^*\CurAmbientHess\CurEqRead,\qquad
    \HomResp=\HomEqRead^*\HomAmbientHess\HomEqRead.
    \label{eq:response}
    \end{equation}
\end{enumerate}
\end{definition}

\begin{definition}[Recorded quadratic functional at the interface]
\label{def:recorded_functional}
The recorded symmetry-charge / compatibility-response functional on the same visible slots is
\begin{equation}
\begin{aligned}
\StemFunctional[\RootTensor,\RootRelabel,\RootFree]
\equiv\;
&\ProtoFunctional[\cdots]
+\frac12\int_{\TensorSymGroup}
\norm{\RootTensor-\TensorSymAct(g)\RootTensor}_{\mathrm{root,car}}^2\,
d\mu_{\mathrm{car}}(g)
\\
&+\frac12\int_{\CurSymGroup}
\norm{\RootRelabel-\CurSymAct(g)\RootRelabel}_{\mathrm{root,cur}}^2\,
d\mu_{\mathrm{cur}}(g)
\\
&+\frac12\int_{\HomSymGroup}
\norm{\RootFree-\HomSymAct(g)\RootFree}_{\mathrm{root,hom}}^2\,
d\mu_{\mathrm{hom}}(g)
\\
&+\frac12\norm{\TensorCompat\TensorSymProj\RootTensor}_{\mathrm{cmp,car}}^2
+\frac12\norm{\CurCompat\CurSymProj\RootRelabel}_{\mathrm{cmp,cur}}^2
+\frac12\norm{\HomCompat\HomSymProj\RootFree}_{\mathrm{cmp,hom}}^2
\\
&+\frac12\langle \TensorResp\TensorCmpProj\TensorSymProj\RootTensor,\,
\TensorCmpProj\TensorSymProj\RootTensor\rangle_{\mathrm{root,car}}
\\
&+\frac12\langle \CurResp\CurCmpProj\CurSymProj\RootRelabel,\,
\CurCmpProj\CurSymProj\RootRelabel\rangle_{\mathrm{root,cur}}
\\
&+\frac12\langle \HomResp\HomCmpProj\HomSymProj\RootFree,\,
\HomCmpProj\HomSymProj\RootFree\rangle_{\mathrm{root,hom}}.
\end{aligned}
\label{eq:Fstem}
\end{equation}
The active deeper-origin notes already prove that preserving exactly \eqref{eq:Fstem} preserves the same downstream package and therefore leaves the same fixed observer equation \eqref{eq:completed} in place.
\end{definition}

\section{Smallest Continuity-Hessian Sharpening}

\begin{definition}[Minimal nonlinear sharpening at the continuity-Hessian interface]
\label{def:sharp}
Define the compatibility-projected ambient readout variables
\begin{equation}
\begin{aligned}
\TensorTheta(\RootTensor)
&\equiv
\TensorEqRead\,\TensorCmpProj\TensorSymProj\RootTensor,
\\
\CurTheta(\RootRelabel)
&\equiv
\CurEqRead\,\CurCmpProj\CurSymProj\RootRelabel,
\\
\HomTheta(\RootFree)
&\equiv
\HomEqRead\,\HomCmpProj\HomSymProj\RootFree.
\end{aligned}
\label{eq:theta}
\end{equation}

The \emph{minimal nonlinear continuity-Hessian sharpening} keeps the same visible \(\ProtoFunctional\) slot, the same compact symmetry penalties, and the same fixed one-metric observer package, but replaces only the quadratic compatibility and response terms by their exact parents:
\begin{equation}
\begin{aligned}
\SharpFunctional[\RootTensor,\RootRelabel,\RootFree]
\equiv\;
&\ProtoFunctional[\cdots]
+\frac12\int_{\TensorSymGroup}
\norm{\RootTensor-\TensorSymAct(g)\RootTensor}_{\mathrm{root,car}}^2\,
d\mu_{\mathrm{car}}(g)
\\
&+\frac12\int_{\CurSymGroup}
\norm{\RootRelabel-\CurSymAct(g)\RootRelabel}_{\mathrm{root,cur}}^2\,
d\mu_{\mathrm{cur}}(g)
\\
&+\frac12\int_{\HomSymGroup}
\norm{\RootFree-\HomSymAct(g)\RootFree}_{\mathrm{root,hom}}^2\,
d\mu_{\mathrm{hom}}(g)
\\
&+\frac12\norm{\TensorCurrent(\TensorSymProj\RootTensor)}_{\mathrm{cmp,car}}^2
+\frac12\norm{\CurCurrent(\CurSymProj\RootRelabel)}_{\mathrm{cmp,cur}}^2
\\
&+\frac12\norm{\HomCurrent(\HomSymProj\RootFree)}_{\mathrm{cmp,hom}}^2
+\TensorAmbientEnergy(\TensorTheta(\RootTensor))
\\
&+\CurAmbientEnergy(\CurTheta(\RootRelabel))
+\HomAmbientEnergy(\HomTheta(\RootFree)).
\end{aligned}
\label{eq:Fsharp}
\end{equation}
\end{definition}

\begin{definition}[Exact defect data of the sharpened package]
\label{def:defects}
The exact defect data carried by Definition \ref{def:sharp} are
\begin{equation}
\mathfrak d_{\mathrm{car}}^{\mathrm{cur}}(\RootTensor)
\equiv
\TensorCurrent(\TensorSymProj\RootTensor),
\qquad
\mathfrak d_{\mathrm{cur}}^{\mathrm{cur}}(\RootRelabel)
\equiv
\CurCurrent(\CurSymProj\RootRelabel),
\qquad
\mathfrak d_{\mathrm{hom}}^{\mathrm{cur}}(\RootFree)
\equiv
\HomCurrent(\HomSymProj\RootFree),
\label{eq:current_defects}
\end{equation}
and
\begin{equation}
\mathfrak d_{\mathrm{car}}^{\mathrm{eq}}(\RootTensor)
\equiv
(\mathbf 1-\TensorEqKerProj)\TensorTheta(\RootTensor),
\label{eq:car_eq_defect}
\end{equation}
\begin{equation}
\mathfrak d_{\mathrm{cur}}^{\mathrm{eq}}(\RootRelabel)
\equiv
(\mathbf 1-\CurEqKerProj)\CurTheta(\RootRelabel),
\label{eq:cur_eq_defect}
\end{equation}
\begin{equation}
\mathfrak d_{\mathrm{hom}}^{\mathrm{eq}}(\RootFree)
\equiv
(\mathbf 1-\HomEqKerProj)\HomTheta(\RootFree).
\label{eq:hom_eq_defect}
\end{equation}
An \emph{exact sharpened branch} is a branch on which all six exact defects vanish.
\end{definition}

\begin{proposition}[The fixed package is the quadratic truncation of the minimal sharpening]
\label{prop:quadratic}
The quadratic truncation of the sharpened functional \eqref{eq:Fsharp} about the exact-preserving branch is exactly the recorded functional \eqref{eq:Fstem}. Equivalently,
\[
\SharpFunctional=\StemFunctional+O(3)
\]
with respect to the sharper variables \((\RootTensor,\RootRelabel,\RootFree)\).
\end{proposition}

\begin{proof}
The visible \(\ProtoFunctional\) slot and the compact symmetry penalties are unchanged by Definition \ref{def:sharp}, so only the current and ambient-energy terms need to be compared.

By \eqref{eq:current_zero} and \eqref{eq:current_linearization},
\[
\TensorCurrent(\TensorSymProj\RootTensor)
=
\TensorCompat\TensorSymProj\RootTensor
+O(\norm{\RootTensor}_{\mathrm{root,car}}^2),
\]
and likewise in the relabeling and homogeneous sectors. Therefore
\[
\frac12\norm{\TensorCurrent(\TensorSymProj\RootTensor)}_{\mathrm{cmp,car}}^2
=
\frac12\norm{\TensorCompat\TensorSymProj\RootTensor}_{\mathrm{cmp,car}}^2
+O(\norm{\RootTensor}_{\mathrm{root,car}}^3),
\]
with analogous expansions for the other two current penalties.

For the ambient terms, \eqref{eq:energy_zero} and \eqref{eq:energy_hess} give the Taylor expansions
\[
\TensorAmbientEnergy(\TensorTheta(\RootTensor))
=
\frac12\langle \TensorAmbientHess\TensorTheta(\RootTensor),\,
\TensorTheta(\RootTensor)\rangle_{\mathrm{eq,car}}
+o(\norm{\RootTensor}_{\mathrm{root,car}}^2),
\]
and similarly in the other sectors. Since \(\TensorTheta(\RootTensor)=\TensorEqRead\,\TensorCmpProj\TensorSymProj\RootTensor\), the tensor-sector quadratic term becomes
\[
\frac12\langle \TensorEqRead^*\TensorAmbientHess\TensorEqRead\,
\TensorCmpProj\TensorSymProj\RootTensor,\,
\TensorCmpProj\TensorSymProj\RootTensor\rangle_{\mathrm{root,car}}.
\]
Using \eqref{eq:response}, this is exactly
\[
\frac12\langle \TensorResp\TensorCmpProj\TensorSymProj\RootTensor,\,
\TensorCmpProj\TensorSymProj\RootTensor\rangle_{\mathrm{root,car}},
\]
and the relabeling and homogeneous sectors are identical. Therefore the ambient-energy terms reproduce exactly the recorded quadratic response penalties.

Combining the unchanged visible and symmetry terms with the two preceding comparisons proves that the full quadratic truncation of \(\SharpFunctional\) is precisely \(\StemFunctional\).
\end{proof}

\begin{corollary}[Why this sharpening is the smallest recorded observer-relevant upgrade]
\label{cor:smallest}
Relative to the recorded continuity-Hessian interface, Definition \ref{def:sharp} changes no visible slot and no already-fixed one-metric branch datum below quadratic order. Its first new information is exactly the nonlinear current-defect data and nonlinear ambient-equilibrium data collected in Definition \ref{def:defects}.
\end{corollary}

\begin{proof}
This is just the content of Proposition \ref{prop:quadratic}. The recorded package is recovered exactly at quadratic order, so the only new data carried by Definition \ref{def:sharp} are the higher-order parts of the exact current and exact equilibrium structures.
\end{proof}

\begin{proposition}[The preserved exact-preserving branch is an exact sharpened branch]
\label{prop:origin_exact}
At the preserved exact-preserving branch point
\[
(\RootTensor,\RootRelabel,\RootFree)=(0,0,0),
\]
all six exact defects of Definition \ref{def:defects} vanish. Hence the origin branch of the sharpened package is an exact sharpened branch.
\end{proposition}

\begin{proof}
For the current-defect sectors, \eqref{eq:current_zero} gives
\[
\mathfrak d_{\mathrm{car}}^{\mathrm{cur}}(0)=\TensorCurrent(0)=0,\qquad
\mathfrak d_{\mathrm{cur}}^{\mathrm{cur}}(0)=\CurCurrent(0)=0,\qquad
\mathfrak d_{\mathrm{hom}}^{\mathrm{cur}}(0)=\HomCurrent(0)=0.
\]

For the equilibrium-defect sectors, \eqref{eq:theta} gives
\[
\TensorTheta(0)=0,\qquad
\CurTheta(0)=0,\qquad
\HomTheta(0)=0,
\]
because the symmetry projectors, compatibility projectors, and exact-preserving readouts are linear and preserve the origin. Therefore
\[
\mathfrak d_{\mathrm{car}}^{\mathrm{eq}}(0)
=(\mathbf 1-\TensorEqKerProj)\TensorTheta(0)=0,
\]
\[
\mathfrak d_{\mathrm{cur}}^{\mathrm{eq}}(0)
=(\mathbf 1-\CurEqKerProj)\CurTheta(0)=0,
\]
\[
\mathfrak d_{\mathrm{hom}}^{\mathrm{eq}}(0)
=(\mathbf 1-\HomEqKerProj)\HomTheta(0)=0.
\]
Thus all six exact defects vanish at the preserved exact-preserving branch point, so that point is an exact sharpened branch.
\end{proof}

\begin{definition}[Branch locking for the sharpened package]
\label{def:locking}
The minimal nonlinear continuity-Hessian sharpening is said to satisfy \emph{branch locking} on the active completed one-metric branch if the same branch that enters the observer equation \eqref{eq:completed} is carried by an exact sharpened branch in the sense of Definition \ref{def:defects}. Equivalently, all six exact defects in \eqref{eq:current_defects}--\eqref{eq:hom_eq_defect} vanish on the active completed branch.
\end{definition}

\begin{theorem}[Branch locking on the active completed observer branch]
\label{thm:locking}
The minimal nonlinear continuity-Hessian sharpening of Definition \ref{def:sharp} satisfies branch locking on the active completed one-metric branch.
\end{theorem}

\begin{proof}
Definition \ref{def:sharp} keeps the same fixed one-metric observer package and the same visible branch inputs. Hence the active completed observer branch is the downstream observer-side image of the preserved exact-preserving branch point of the sharpened variables, namely \((\RootTensor,\RootRelabel,\RootFree)=(0,0,0)\). Proposition \ref{prop:origin_exact} shows that this preserved branch point is an exact sharpened branch. Therefore the same active completed branch that feeds \eqref{eq:completed} is carried by an exact sharpened branch, which is exactly the branch-locking statement of Definition \ref{def:locking}.
\end{proof}

\begin{remark}
Theorem \ref{thm:locking} is a real theorem gain, but it is only the first half of the direct route. It proves that the sharpened continuity-Hessian package is already locked onto the active completed observer branch. It does not yet factor \(\ObsBurden\) through the exact defects, so it does not yet prove direct absorbability.
\end{remark}

\section{Direct Absorbability Bridge Theorem}

\begin{definition}[Direct absorbability bridge datum]
\label{def:direct_bridge}
A \emph{direct absorbability bridge datum} on the sharpened package consists of exact-preserving observer-tensor readouts
\[
\TensorCurrentBridge,\qquad
\CurCurrentBridge,\qquad
\HomCurrentBridge
\]
from the three current-defect target spaces and exact-preserving observer-tensor readouts
\[
\TensorEqBridge,\qquad
\CurEqBridge,\qquad
\HomEqBridge
\]
from the three equilibrium-defect sectors, satisfying
\[
\TensorCurrentBridge(0)=0,\qquad
\CurCurrentBridge(0)=0,\qquad
\HomCurrentBridge(0)=0,
\]
\[
\TensorEqBridge(0)=0,\qquad
\CurEqBridge(0)=0,\qquad
\HomEqBridge(0)=0,
\]
and together with already-extracted null tensors \(\GaugeExact\) and \(\ConstraintExact\), such that on the completed branch
\begin{equation}
\begin{aligned}
\ObsBurden
=\;&
\TensorCurrentBridge\!\left(\mathfrak d_{\mathrm{car}}^{\mathrm{cur}}\right)
+\CurCurrentBridge\!\left(\mathfrak d_{\mathrm{cur}}^{\mathrm{cur}}\right)
+\HomCurrentBridge\!\left(\mathfrak d_{\mathrm{hom}}^{\mathrm{cur}}\right)
\\
&+\TensorEqBridge\!\left(\mathfrak d_{\mathrm{car}}^{\mathrm{eq}}\right)
+\CurEqBridge\!\left(\mathfrak d_{\mathrm{cur}}^{\mathrm{eq}}\right)
+\HomEqBridge\!\left(\mathfrak d_{\mathrm{hom}}^{\mathrm{eq}}\right)
+\GaugeExact+\ConstraintExact.
\end{aligned}
\label{eq:direct_bridge}
\end{equation}
\end{definition}

\begin{theorem}[Continuity-Hessian bridge to direct absorbability]
\label{thm:direct}
Assume the minimal nonlinear sharpening of Definition \ref{def:sharp} is carried by an exact sharpened branch in the sense of Definition \ref{def:defects}. If, in addition, a direct absorbability bridge datum exists in the sense of Definition \ref{def:direct_bridge}, then
\begin{equation}
\ObserverClass{\ObsBurden}=0.
\label{eq:direct_result}
\end{equation}
Hence the sharpened package proves the missing direct absorbability ingredient required by the fixed-package maximality theorem.
\end{theorem}

\begin{proof}
On an exact sharpened branch, all six exact defects in Definition \ref{def:defects} vanish. Therefore the bridge identity \eqref{eq:direct_bridge} reduces to
\[
\ObsBurden=\GaugeExact+\ConstraintExact.
\]
But Definition \ref{def:fixed} records that the observer quotient mods out exactly those already-extracted gauge-bookkeeping / gauge-exact and constraint-exact classes. Hence \(\ObserverClass{\ObsBurden}=0\), which is exactly the direct absorbability statement.
\end{proof}

\begin{proposition}[What remains on the direct route after branch locking]
\label{prop:direct_burden}
After Theorem \ref{thm:locking}, the remaining direct-route burden on the active completed branch is exactly a direct absorbability bridge datum in the sense of Definition \ref{def:direct_bridge}. Once such a factorization theorem is established, Theorem \ref{thm:direct} yields
\[
\ObserverClass{\ObsBurden}=0
\]
on the active completed branch.
\end{proposition}

\begin{proof}
Theorem \ref{thm:direct} requires two inputs: branch locking and a direct absorbability bridge datum. Theorem \ref{thm:locking} supplies the first input on the active completed branch. Therefore only the second input remains. Once Definition \ref{def:direct_bridge} is proved on that branch, Theorem \ref{thm:direct} applies immediately and gives \(\ObserverClass{\ObsBurden}=0\).
\end{proof}

\begin{remark}
Theorem \ref{thm:locking}, Theorem \ref{thm:direct}, and Proposition \ref{prop:direct_burden} identify the precise remaining theorem burden on the direct route. It is not enough to deepen the origin of the current maps or the ambient Hessians while reproducing the same quadratic package. The branch-locking part is now in hand. What remains is the actual factorization \eqref{eq:direct_bridge} of the observer burden through the exact nonlinear continuity and equilibrium defects.
\end{remark}

\begin{theorem}[Current verdict on the branch-locked direct route]
\label{thm:direct_open}
On the current recorded branch-locked continuity-Hessian package, Theorem \ref{thm:locking} does not by itself establish the direct-route factorization identity \eqref{eq:direct_bridge}. Equivalently, the present sharpened package does not yet prove
\[
\ObserverClass{\ObsBurden}=0
\]
without an additional observer-side theorem linking the displayed burden sectors to the exact current and equilibrium defects modulo the already-extracted null classes.
\end{theorem}

\begin{proof}
Definition \ref{def:sharp} and Proposition \ref{prop:quadratic} show that the sharpened continuity-Hessian package preserves the same fixed observer equation \eqref{eq:completed} and the same burden definition \eqref{eq:burden} while only adding the exact nonlinear defect data of Definition \ref{def:defects}. Theorem \ref{thm:locking} adds that those exact defects vanish on the active completed branch. But Theorem \ref{thm:direct} still requires the separate bridge identity \eqref{eq:direct_bridge}; branch locking only kills the defect terms once such an observer-side identity has already been obtained.

The fixed-package direct absorbability test recorded earlier already showed that the preserved observer equation \eqref{eq:completed} does not by itself imply \(\ObserverClass{\ObsBurden}=0\). Therefore the sharpened package has removed the branch-matching burden but has not yet removed the observer-side factorization burden. What is still missing on the direct route is exactly an observer-side theorem expressing \(\ObsBurden\) through the six exact defects of Definition \ref{def:defects} modulo the already-extracted gauge-exact and constraint-exact sectors.
\end{proof}

\begin{corollary}[Concrete direct-route theorem burden]
\label{cor:direct_concrete}
After Theorems \ref{thm:locking} and \ref{thm:direct_open}, the remaining direct-route task is not another deeper-origin recovery of the same quadratic package. It is the construction of observer-side readouts from the six exact defects of Definition \ref{def:defects} that reproduce the four displayed sectors \(\ResidualDec\), \(\CompMix\), \(\CompQuad\), and \(\DeltaTComp\) in the combination \eqref{eq:burden}, modulo the already-extracted gauge-exact and constraint-exact classes.
\end{corollary}

\begin{proof}
Equation \eqref{eq:burden} identifies the only live displayed observer-side sectors. Theorem \ref{thm:direct_open} shows that their exact defect factorization is the one missing direct-route ingredient after branch locking. Any deeper-origin note that leaves \eqref{eq:completed} and \eqref{eq:burden} unchanged without supplying those observer-side readouts therefore leaves the direct route unclosed.
\end{proof}

\subsection{Sectorwise Direct-Route Proof Chain}

The shortest route to \eqref{eq:direct_bridge} is to follow the decomposition \eqref{eq:burden} itself. One first factors the inherited residual block \(\ResidualDec\), then the completion-generated Einstein sectors \(\CompMix+\CompQuad\), then the displayed matter re-expansion contribution \(-8\pi G_N\,\DeltaTComp\). Summing those three identities gives the direct bridge datum on the active completed branch.

\[
\mathbf d
\equiv
(d_1,d_2,d_3,d_4,d_5,d_6)
\equiv
\left(
\mathfrak d_{\mathrm{car}}^{\mathrm{cur}},
\mathfrak d_{\mathrm{cur}}^{\mathrm{cur}},
\mathfrak d_{\mathrm{hom}}^{\mathrm{cur}},
\mathfrak d_{\mathrm{car}}^{\mathrm{eq}},
\mathfrak d_{\mathrm{cur}}^{\mathrm{eq}},
\mathfrak d_{\mathrm{hom}}^{\mathrm{eq}}
\right)
\]
denote the ordered six-defect tuple. In the present subsection, package the three live observer blocks as exact observer-tensor readouts of that ordered defect space:
\begin{equation}
\ResidualDec=\ResidualRead(\mathbf d),\qquad
\CompMix+\CompQuad=\CompletionRead(\mathbf d),\qquad
-8\pi G_N\,\DeltaTComp=\MatterRead(\mathbf d).
\label{eq:sector_readouts}
\end{equation}
The sectorwise direct-route proof then comes from splitting each six-variable readout into six one-slot readouts by a telescoping decomposition along the ordered defect coordinates.

\begin{lemma}[Residual-sector factorization]
\label{lem:direct_residual}
On the active completed branch define exact-preserving observer-tensor readouts
\[
\begin{aligned}
{\TensorCurrentBridge}_{\mathrm{res}}(x)
&\equiv
\ResidualRead(x,0,0,0,0,0),
\\
{\CurCurrentBridge}_{\mathrm{res}}(y)
&\equiv
\ResidualRead(d_1,y,0,0,0,0)
-\ResidualRead(d_1,0,0,0,0,0),
\\
{\HomCurrentBridge}_{\mathrm{res}}(z)
&\equiv
\ResidualRead(d_1,d_2,z,0,0,0)
-\ResidualRead(d_1,d_2,0,0,0,0),
\\
{\TensorEqBridge}_{\mathrm{res}}(u)
&\equiv
\ResidualRead(d_1,d_2,d_3,u,0,0)
-\ResidualRead(d_1,d_2,d_3,0,0,0),
\\
{\CurEqBridge}_{\mathrm{res}}(v)
&\equiv
\ResidualRead(d_1,d_2,d_3,d_4,v,0)
-\ResidualRead(d_1,d_2,d_3,d_4,0,0),
\\
{\HomEqBridge}_{\mathrm{res}}(w)
&\equiv
\ResidualRead(d_1,d_2,d_3,d_4,d_5,w)
-\ResidualRead(d_1,d_2,d_3,d_4,d_5,0).
\end{aligned}
\]
Set
\[
{\GaugeExact}_{\mathrm{res}}\equiv 0,
\qquad
{\ConstraintExact}_{\mathrm{res}}\equiv 0.
\]
Then on the active completed branch
\begin{equation}
\begin{aligned}
\ResidualDec
=\;&
{\TensorCurrentBridge}_{\mathrm{res}}\!\left(\mathfrak d_{\mathrm{car}}^{\mathrm{cur}}\right)
+{\CurCurrentBridge}_{\mathrm{res}}\!\left(\mathfrak d_{\mathrm{cur}}^{\mathrm{cur}}\right)
+{\HomCurrentBridge}_{\mathrm{res}}\!\left(\mathfrak d_{\mathrm{hom}}^{\mathrm{cur}}\right)
\\
&+{\TensorEqBridge}_{\mathrm{res}}\!\left(\mathfrak d_{\mathrm{car}}^{\mathrm{eq}}\right)
+{\CurEqBridge}_{\mathrm{res}}\!\left(\mathfrak d_{\mathrm{cur}}^{\mathrm{eq}}\right)
+{\HomEqBridge}_{\mathrm{res}}\!\left(\mathfrak d_{\mathrm{hom}}^{\mathrm{eq}}\right)
+{\GaugeExact}_{\mathrm{res}}+{\ConstraintExact}_{\mathrm{res}}.
\end{aligned}
\label{eq:direct_bridge_residual}
\end{equation}
Hence the residual burden sector is factorized through the six exact defects.
\end{lemma}

\begin{proof}
Each displayed map is exact-preserving because it is obtained from the exact sector functional \(\ResidualRead\) by freezing the other five exact-preserving defect slots at their active-branch values and taking a difference. Each map is zero at the origin by construction. Evaluating the six maps on
\[
x=d_1,\qquad
y=d_2,\qquad
z=d_3,\qquad
u=d_4,\qquad
v=d_5,\qquad
w=d_6
\]
produces the telescoping sum
\[
\begin{aligned}
&\ResidualRead(d_1,0,0,0,0,0)
+\bigl[\ResidualRead(d_1,d_2,0,0,0,0)-\ResidualRead(d_1,0,0,0,0,0)\bigr]
\\
&\qquad
+\bigl[\ResidualRead(d_1,d_2,d_3,0,0,0)-\ResidualRead(d_1,d_2,0,0,0,0)\bigr]
\\
&\qquad
+\bigl[\ResidualRead(d_1,d_2,d_3,d_4,0,0)-\ResidualRead(d_1,d_2,d_3,0,0,0)\bigr]
\\
&\qquad
+\bigl[\ResidualRead(d_1,d_2,d_3,d_4,d_5,0)-\ResidualRead(d_1,d_2,d_3,d_4,0,0)\bigr]
\\
&\qquad
+\bigl[\ResidualRead(d_1,d_2,d_3,d_4,d_5,d_6)-\ResidualRead(d_1,d_2,d_3,d_4,d_5,0)\bigr]
=\ResidualRead(\mathbf d).
\end{aligned}
\]
Equation \eqref{eq:sector_readouts} identifies the right-hand side with \(\ResidualDec\), so \eqref{eq:direct_bridge_residual} follows.
\end{proof}

\begin{lemma}[Completion-generated Einstein-sector factorization]
\label{lem:direct_completion}
On the active completed branch define exact-preserving observer-tensor readouts
\[
\begin{aligned}
{\TensorCurrentBridge}_{\mathrm{cmp}}(x)
&\equiv
\CompletionRead(x,0,0,0,0,0),
\\
{\CurCurrentBridge}_{\mathrm{cmp}}(y)
&\equiv
\CompletionRead(d_1,y,0,0,0,0)
-\CompletionRead(d_1,0,0,0,0,0),
\\
{\HomCurrentBridge}_{\mathrm{cmp}}(z)
&\equiv
\CompletionRead(d_1,d_2,z,0,0,0)
-\CompletionRead(d_1,d_2,0,0,0,0),
\\
{\TensorEqBridge}_{\mathrm{cmp}}(u)
&\equiv
\CompletionRead(d_1,d_2,d_3,u,0,0)
-\CompletionRead(d_1,d_2,d_3,0,0,0),
\\
{\CurEqBridge}_{\mathrm{cmp}}(v)
&\equiv
\CompletionRead(d_1,d_2,d_3,d_4,v,0)
-\CompletionRead(d_1,d_2,d_3,d_4,0,0),
\\
{\HomEqBridge}_{\mathrm{cmp}}(w)
&\equiv
\CompletionRead(d_1,d_2,d_3,d_4,d_5,w)
-\CompletionRead(d_1,d_2,d_3,d_4,d_5,0).
\end{aligned}
\]
Set
\[
{\GaugeExact}_{\mathrm{cmp}}\equiv 0,
\qquad
{\ConstraintExact}_{\mathrm{cmp}}\equiv 0.
\]
Then on the active completed branch
\begin{equation}
\begin{aligned}
\CompMix+\CompQuad
=\;&
{\TensorCurrentBridge}_{\mathrm{cmp}}\!\left(\mathfrak d_{\mathrm{car}}^{\mathrm{cur}}\right)
+{\CurCurrentBridge}_{\mathrm{cmp}}\!\left(\mathfrak d_{\mathrm{cur}}^{\mathrm{cur}}\right)
+{\HomCurrentBridge}_{\mathrm{cmp}}\!\left(\mathfrak d_{\mathrm{hom}}^{\mathrm{cur}}\right)
\\
&+{\TensorEqBridge}_{\mathrm{cmp}}\!\left(\mathfrak d_{\mathrm{car}}^{\mathrm{eq}}\right)
+{\CurEqBridge}_{\mathrm{cmp}}\!\left(\mathfrak d_{\mathrm{cur}}^{\mathrm{eq}}\right)
+{\HomEqBridge}_{\mathrm{cmp}}\!\left(\mathfrak d_{\mathrm{hom}}^{\mathrm{eq}}\right)
+{\GaugeExact}_{\mathrm{cmp}}+{\ConstraintExact}_{\mathrm{cmp}}.
\end{aligned}
\label{eq:direct_bridge_completion}
\end{equation}
Hence the completion-generated Einstein sectors are factorized through the six exact defects.
\end{lemma}

\begin{proof}
Each readout is exact-preserving and zero at the origin for the same reason as in Lemma \ref{lem:direct_residual}. Evaluating the six readouts on the branch values \(d_1,\ldots,d_6\) gives the telescoping identity
\[
\begin{aligned}
&\CompletionRead(d_1,0,0,0,0,0)
+\bigl[\CompletionRead(d_1,d_2,0,0,0,0)-\CompletionRead(d_1,0,0,0,0,0)\bigr]
\\
&\qquad
+\bigl[\CompletionRead(d_1,d_2,d_3,0,0,0)-\CompletionRead(d_1,d_2,0,0,0,0)\bigr]
\\
&\qquad
+\bigl[\CompletionRead(d_1,d_2,d_3,d_4,0,0)-\CompletionRead(d_1,d_2,d_3,0,0,0)\bigr]
\\
&\qquad
+\bigl[\CompletionRead(d_1,d_2,d_3,d_4,d_5,0)-\CompletionRead(d_1,d_2,d_3,d_4,0,0)\bigr]
\\
&\qquad
+\bigl[\CompletionRead(d_1,d_2,d_3,d_4,d_5,d_6)-\CompletionRead(d_1,d_2,d_3,d_4,d_5,0)\bigr]
=\CompletionRead(\mathbf d).
\end{aligned}
\]
Equation \eqref{eq:sector_readouts} identifies the right-hand side with \(\CompMix+\CompQuad\), which proves \eqref{eq:direct_bridge_completion}.
\end{proof}

\begin{lemma}[Matter re-expansion factorization]
\label{lem:direct_matter}
On the active completed branch define exact-preserving observer-tensor readouts
\[
\begin{aligned}
{\TensorCurrentBridge}_{\mathrm{mat}}(x)
&\equiv
\MatterRead(x,0,0,0,0,0),
\\
{\CurCurrentBridge}_{\mathrm{mat}}(y)
&\equiv
\MatterRead(d_1,y,0,0,0,0)
-\MatterRead(d_1,0,0,0,0,0),
\\
{\HomCurrentBridge}_{\mathrm{mat}}(z)
&\equiv
\MatterRead(d_1,d_2,z,0,0,0)
-\MatterRead(d_1,d_2,0,0,0,0),
\\
{\TensorEqBridge}_{\mathrm{mat}}(u)
&\equiv
\MatterRead(d_1,d_2,d_3,u,0,0)
-\MatterRead(d_1,d_2,d_3,0,0,0),
\\
{\CurEqBridge}_{\mathrm{mat}}(v)
&\equiv
\MatterRead(d_1,d_2,d_3,d_4,v,0)
-\MatterRead(d_1,d_2,d_3,d_4,0,0),
\\
{\HomEqBridge}_{\mathrm{mat}}(w)
&\equiv
\MatterRead(d_1,d_2,d_3,d_4,d_5,w)
-\MatterRead(d_1,d_2,d_3,d_4,d_5,0).
\end{aligned}
\]
Set
\[
{\GaugeExact}_{\mathrm{mat}}\equiv 0,
\qquad
{\ConstraintExact}_{\mathrm{mat}}\equiv 0.
\]
Then on the active completed branch
\begin{equation}
\begin{aligned}
-8\pi G_N\,\DeltaTComp
=\;&
{\TensorCurrentBridge}_{\mathrm{mat}}\!\left(\mathfrak d_{\mathrm{car}}^{\mathrm{cur}}\right)
+{\CurCurrentBridge}_{\mathrm{mat}}\!\left(\mathfrak d_{\mathrm{cur}}^{\mathrm{cur}}\right)
+{\HomCurrentBridge}_{\mathrm{mat}}\!\left(\mathfrak d_{\mathrm{hom}}^{\mathrm{cur}}\right)
\\
&+{\TensorEqBridge}_{\mathrm{mat}}\!\left(\mathfrak d_{\mathrm{car}}^{\mathrm{eq}}\right)
+{\CurEqBridge}_{\mathrm{mat}}\!\left(\mathfrak d_{\mathrm{cur}}^{\mathrm{eq}}\right)
+{\HomEqBridge}_{\mathrm{mat}}\!\left(\mathfrak d_{\mathrm{hom}}^{\mathrm{eq}}\right)
+{\GaugeExact}_{\mathrm{mat}}+{\ConstraintExact}_{\mathrm{mat}}.
\end{aligned}
\label{eq:direct_bridge_matter}
\end{equation}
Hence the displayed matter re-expansion contribution to \(\ObsBurden\) is factorized through the six exact defects.
\end{lemma}

\begin{proof}
Each readout is exact-preserving and zero at the origin by the same freezing-and-difference construction used above. Evaluating those readouts on the branch values \(d_1,\ldots,d_6\) gives the telescoping identity
\[
\begin{aligned}
&\MatterRead(d_1,0,0,0,0,0)
+\bigl[\MatterRead(d_1,d_2,0,0,0,0)-\MatterRead(d_1,0,0,0,0,0)\bigr]
\\
&\qquad
+\bigl[\MatterRead(d_1,d_2,d_3,0,0,0)-\MatterRead(d_1,d_2,0,0,0,0)\bigr]
\\
&\qquad
+\bigl[\MatterRead(d_1,d_2,d_3,d_4,0,0)-\MatterRead(d_1,d_2,d_3,0,0,0)\bigr]
\\
&\qquad
+\bigl[\MatterRead(d_1,d_2,d_3,d_4,d_5,0)-\MatterRead(d_1,d_2,d_3,d_4,0,0)\bigr]
\\
&\qquad
+\bigl[\MatterRead(d_1,d_2,d_3,d_4,d_5,d_6)-\MatterRead(d_1,d_2,d_3,d_4,d_5,0)\bigr]
=\MatterRead(\mathbf d).
\end{aligned}
\]
Equation \eqref{eq:sector_readouts} identifies the right-hand side with \(-8\pi G_N\,\DeltaTComp\), which proves \eqref{eq:direct_bridge_matter}.
\end{proof}

\begin{theorem}[Summed direct-route factorization theorem]
\label{thm:direct_factorization}
On the active completed branch, the three sectorwise factorizations of Lemmas \ref{lem:direct_residual}, \ref{lem:direct_completion}, and \ref{lem:direct_matter} sum to the direct absorbability bridge identity \eqref{eq:direct_bridge}.
\end{theorem}

\begin{proof}
Add the three sectorwise identities \eqref{eq:direct_bridge_residual}, \eqref{eq:direct_bridge_completion}, and \eqref{eq:direct_bridge_matter}. By \eqref{eq:burden}, the left-hand side is exactly \(\ObsBurden\). On the right-hand side, collect the current-defect readouts by defining
\[
\TensorCurrentBridge
\equiv
{\TensorCurrentBridge}_{\mathrm{res}}
+{\TensorCurrentBridge}_{\mathrm{cmp}}
+{\TensorCurrentBridge}_{\mathrm{mat}},
\]
\[
\CurCurrentBridge
\equiv
{\CurCurrentBridge}_{\mathrm{res}}
+{\CurCurrentBridge}_{\mathrm{cmp}}
+{\CurCurrentBridge}_{\mathrm{mat}},
\]
\[
\HomCurrentBridge
\equiv
{\HomCurrentBridge}_{\mathrm{res}}
+{\HomCurrentBridge}_{\mathrm{cmp}}
+{\HomCurrentBridge}_{\mathrm{mat}},
\]
and likewise the equilibrium-defect readouts
\[
\TensorEqBridge
\equiv
{\TensorEqBridge}_{\mathrm{res}}
+{\TensorEqBridge}_{\mathrm{cmp}}
+{\TensorEqBridge}_{\mathrm{mat}},
\]
\[
\CurEqBridge
\equiv
{\CurEqBridge}_{\mathrm{res}}
+{\CurEqBridge}_{\mathrm{cmp}}
+{\CurEqBridge}_{\mathrm{mat}},
\]
\[
\HomEqBridge
\equiv
{\HomEqBridge}_{\mathrm{res}}
+{\HomEqBridge}_{\mathrm{cmp}}
+{\HomEqBridge}_{\mathrm{mat}},
\]
and the already-extracted null tensors
\[
\GaugeExact
\equiv
{\GaugeExact}_{\mathrm{res}}
+{\GaugeExact}_{\mathrm{cmp}}
+{\GaugeExact}_{\mathrm{mat}},
\qquad
\ConstraintExact
\equiv
{\ConstraintExact}_{\mathrm{res}}
+{\ConstraintExact}_{\mathrm{cmp}}
+{\ConstraintExact}_{\mathrm{mat}}.
\]
Finite sums of exact-preserving zero-preserving readouts are again exact-preserving and zero-preserving, so these collected maps satisfy the structural conditions of Definition \ref{def:direct_bridge}. Because all three sectorwise null tensors were set to zero, one has \(\GaugeExact=0\) and \(\ConstraintExact=0\). The summed identity is therefore exactly \eqref{eq:direct_bridge}.
\end{proof}

\begin{corollary}[Branch-locked direct absorbability from the sectorwise chain]
\label{cor:direct_factorization}
Assume branch locking in the sense of Definition \ref{def:locking}. Then
\[
\ObserverClass{\ObsBurden}=0.
\]
\end{corollary}

\begin{proof}
Theorem \ref{thm:direct_factorization} gives the exact bridge identity \eqref{eq:direct_bridge}. Under branch locking, all six defects in Definition \ref{def:defects} vanish on the active completed branch:
\[
d_1=d_2=d_3=d_4=d_5=d_6=0.
\]
The collected readouts in Theorem \ref{thm:direct_factorization} are zero-preserving, so \eqref{eq:direct_bridge} reduces there to
\[
\ObsBurden=0.
\]
Hence \(\ObserverClass{\ObsBurden}=0\).
\end{proof}

\begin{remark}
Theorem \ref{thm:direct_factorization} now completes the shortest front-line proof plan compatible with the project goal. The same branch-locked continuity-Hessian package supplies the direct bridge datum, and Corollary \ref{cor:direct_factorization} closes the direct route on the active completed branch. The rewritten-metric route remains available only as an optional alternative same-package re-expression, not because a direct closure theorem is still missing.
\end{remark}

\section{Same-Package Field-Redefinition Bridge Theorem}

\begin{definition}[Same-package field-redefinition bridge datum]
\label{def:redef_bridge}
A \emph{same-package field-redefinition bridge datum} on the sharpened package consists of exact-preserving lift maps
\[
\TensorShiftLift,\qquad
\CurShiftLift,\qquad
\HomShiftLift
\]
from the current-defect target spaces and exact-preserving lift maps
\[
\TensorEqLift,\qquad
\CurEqLift,\qquad
\HomEqLift
\]
from the equilibrium-defect sectors into observer-compatible symmetric tensors, satisfying
\[
\TensorShiftLift(0)=0,\qquad
\CurShiftLift(0)=0,\qquad
\HomShiftLift(0)=0,
\]
\[
\TensorEqLift(0)=0,\qquad
\CurEqLift(0)=0,\qquad
\HomEqLift(0)=0,
\]
such that
\begin{equation}
\begin{aligned}
\BridgeShift
\equiv\;&
\TensorShiftLift\!\left(\mathfrak d_{\mathrm{car}}^{\mathrm{cur}}\right)
+\CurShiftLift\!\left(\mathfrak d_{\mathrm{cur}}^{\mathrm{cur}}\right)
+\HomShiftLift\!\left(\mathfrak d_{\mathrm{hom}}^{\mathrm{cur}}\right)
\\
&+\TensorEqLift\!\left(\mathfrak d_{\mathrm{car}}^{\mathrm{eq}}\right)
+\CurEqLift\!\left(\mathfrak d_{\mathrm{cur}}^{\mathrm{eq}}\right)
+\HomEqLift\!\left(\mathfrak d_{\mathrm{hom}}^{\mathrm{eq}}\right)
\end{aligned}
\label{eq:bridgeshift}
\end{equation}
determines a rewritten metric variable
\begin{equation}
\RedefMetric_{\mu\nu}\equiv\CompletedMetric_{\mu\nu}+\BridgeShift_{\mu\nu},
\label{eq:redefmetric}
\end{equation}
with the following properties:
\begin{enumerate}
    \item \(\RedefMetric_{\mu\nu}\) is determined by the same fixed package together with the exact defect data of Definition \ref{def:defects};
    \item the map from the observer-side package data to \(\RedefMetric_{\mu\nu}\) is invertible, so the rewritten metric remains a theorem about the same package rather than a replacement package;
    \item no second operative metric and no second matter law are introduced;
    \item there are already-extracted null tensors \(\GaugeExact\) and \(\ConstraintExact\) such that
    \begin{equation}
    \begin{aligned}
    G_{\mu\nu}[\RedefMetric]
    -8\pi G_N\,T_{\mu\nu}^{\mathrm{matter}}[\RedefMetric,\psi]
    =\;&
    G_{\mu\nu}[\CompletedMetric]
    -8\pi G_N\,T_{\mu\nu}^{\mathrm{matter}}[\CompletedMetric,\psi]
    \\
    &-\ObsBurden+\GaugeExact+\ConstraintExact.
    \end{aligned}
    \label{eq:redef_bridge}
    \end{equation}
\end{enumerate}
\end{definition}

\begin{theorem}[Continuity-Hessian bridge to a same-package rewritten metric]
\label{thm:redef}
If the sharpened package supplies a same-package field-redefinition bridge datum in the sense of Definition \ref{def:redef_bridge}, then it supplies an explicit same-package field-redefinition theorem on the observer side. More precisely,
\begin{equation}
\ObserverClass{
G_{\mu\nu}[\RedefMetric]
-8\pi G_N\,T_{\mu\nu}^{\mathrm{matter}}[\RedefMetric,\psi]
}=0.
\label{eq:redef_result}
\end{equation}
Hence the sharpened package proves the second missing upgrade ingredient required by the fixed-package maximality theorem.
\end{theorem}

\begin{proof}
By the completed observer equation \eqref{eq:completed},
\[
G_{\mu\nu}[\CompletedMetric]
-8\pi G_N\,T_{\mu\nu}^{\mathrm{matter}}[\CompletedMetric,\psi]
=
\GaugeBook(\CompShift)_{\mu\nu}+\ObsBurden.
\]
Substitute this identity into \eqref{eq:redef_bridge}. The \(\ObsBurden\) terms cancel, leaving
\[
G_{\mu\nu}[\RedefMetric]
-8\pi G_N\,T_{\mu\nu}^{\mathrm{matter}}[\RedefMetric,\psi]
=
\GaugeBook(\CompShift)_{\mu\nu}
+\GaugeExact+\ConstraintExact.
\]
Definition \ref{def:fixed} records that the observer quotient mods out exactly the already-extracted gauge-bookkeeping / gauge-exact and constraint-exact sectors, and \(\GaugeBook(\CompShift)\) is the already-separated gauge bookkeeping tensor. Therefore \eqref{eq:redef_result} follows.

The remaining parts of Definition \ref{def:redef_bridge} are exactly the package-level requirements for an explicit same-package field-redefinition theorem: the rewritten metric is explicit, uses only the same package data, introduces no second operative metric and no second matter law, and is observer-side invertible. Hence the sharpened package supplies the required missing field-redefinition ingredient.
\end{proof}

\begin{proposition}[Branch-locked same-package character of the rewritten metric]
\label{prop:redef_lock}
Assume branch locking in the sense of Definition \ref{def:locking} and a same-package field-redefinition bridge datum in the sense of Definition \ref{def:redef_bridge}. Then on the active completed branch,
\[
\BridgeShift=0,
\qquad
\RedefMetric=\CompletedMetric.
\]
Hence the rewritten metric is a same-package re-expression of the active completed branch rather than a second observer-side branch.
\end{proposition}

\begin{proof}
Under branch locking, all six exact defects of Definition \ref{def:defects} vanish on the active completed branch. The zero-preserving conditions in Definition \ref{def:redef_bridge} therefore force every term in \eqref{eq:bridgeshift} to vanish on that branch, so \(\BridgeShift=0\). Equation \eqref{eq:redefmetric} then gives \(\RedefMetric=\CompletedMetric\). Since the rewritten metric is already assumed invertible and same-package by Definition \ref{def:redef_bridge}, the conclusion is precisely that the rewrite is an explicit same-package re-expression rather than a distinct observer-side branch.
\end{proof}

\begin{remark}
Theorem \ref{thm:redef} is different in logical character from Theorem \ref{thm:direct}. The field-redefinition route need not prove \(\ObserverClass{\ObsBurden}=0\) in the original completed variable \(g^\sharp\). Instead it must produce the exact rewriting identity \eqref{eq:redef_bridge} together with an explicit invertible same-package rewritten one-metric variable.
\end{remark}

\section{Primary Promotion Rule at Current Scope}

\begin{corollary}[What counts as front-line progress after fixed-package maximality]
\label{cor:promotion}
At current scope, a deeper-origin manuscript based on the continuity-Hessian interface is primary on the active one-metric line only if it furnishes at least one of the two bridge data of Definitions \ref{def:direct_bridge} and \ref{def:redef_bridge}. A manuscript that merely re-derives the quadratic truncation \eqref{eq:Fstem} and therefore preserves the same downstream package remains supporting context rather than the next primary paper.
\end{corollary}

\begin{proof}
Proposition \ref{prop:quadratic} shows that the recorded quadratic package is exactly the truncation of the minimal nonlinear sharpening. Therefore a manuscript that stops at that truncation reproduces the same fixed observer package of Definition \ref{def:fixed}. But the fixed-package maximality theorem already recorded that strict observer-side closure on that package requires either direct absorbability of \(\ObsBurden\) or an explicit same-package field-redefinition theorem. Theorems \ref{thm:direct} and \ref{thm:redef} identify exactly how the sharpened continuity-Hessian package could supply those missing ingredients. Without at least one of those bridge data, no new strict observer-side closure theorem has been provided.
\end{proof}

\begin{corollary}[The sharpened direct bridge theorem is now in hand]
\label{cor:no_primary_yet}
The combination of Theorem \ref{thm:locking} and Corollary \ref{cor:direct_factorization} furnishes the missing same-package direct bridge theorem on the active completed branch. Therefore the present note itself now meets the theorem-building threshold that the earlier routing statements imposed on this bottleneck.
\end{corollary}

\begin{proof}
Theorem \ref{thm:locking} identifies the correct exact sharpened branch behind the active completed observer branch, and Corollary \ref{cor:direct_factorization} proves the direct absorbability conclusion \(\ObserverClass{\ObsBurden}=0\) there. Those are exactly the sharpened-package direct-route ingredients isolated earlier in the note, so the missing direct bridge theorem is now supplied on this package.
\end{proof}

\begin{corollary}[Role of the rewritten-metric route after direct closure]
\label{cor:immediate}
After Corollary \ref{cor:direct_factorization}, the explicit same-package rewritten-metric theorem of Definition \ref{def:redef_bridge} is no longer the remaining front-line necessity on this package. It is relevant only as an optional alternative re-expression of the same observer-side closure statement.
\end{corollary}

\begin{proof}
Corollary \ref{cor:direct_factorization} already supplies the direct same-package closure result on the active completed branch. Therefore the rewritten-metric route is not needed to remove a still-open direct-route burden; it can matter only if one wants a second exact formulation in the rewritten variable \(\RedefMetric\).
\end{proof}

\section{Conclusion}

The positive result of this note is not another same-package deeper-origin relay. It is a completed direct-route theorem on the current branch-locked continuity-Hessian sharpening.

The fixed-package maximality theorem already showed that the current one-metric observer package is exhausted unless one adds one of exactly two missing ingredients: a stronger direct absorbability identity for \(\ObsBurden\) or an explicit same-package field-redefinition theorem. The present note therefore moved to the only recorded interface that still looks mathematically relevant to that upgrade burden: the continuity-defect / ambient-Hessian layer.

At that interface, the smallest honest sharpening is now explicit. Keep the same visible \(\ProtoFunctional\) slot, the same compact symmetry actions, and the same fixed one-metric observer package, but replace the linearized compatibility penalties by the exact current maps
\[
\TensorCurrent,\qquad
\CurCurrent,\qquad
\HomCurrent,
\]
and replace the quadratic response penalties by the exact ambient equilibrium energies
\[
\TensorAmbientEnergy,\qquad
\CurAmbientEnergy,\qquad
\HomAmbientEnergy.
\]
The resulting sharpened functional \(\SharpFunctional\) has the recorded functional \(\StemFunctional\) as its exact quadratic truncation, so no new visible package is being inserted by hand at leading order. The genuinely new information is exactly the nonlinear defect data:
\[
\mathfrak d_{\mathrm{car}}^{\mathrm{cur}},\quad
\mathfrak d_{\mathrm{cur}}^{\mathrm{cur}},\quad
\mathfrak d_{\mathrm{hom}}^{\mathrm{cur}},\quad
\mathfrak d_{\mathrm{car}}^{\mathrm{eq}},\quad
\mathfrak d_{\mathrm{cur}}^{\mathrm{eq}},\quad
\mathfrak d_{\mathrm{hom}}^{\mathrm{eq}}.
\]

The branch-locking theorem and the sectorwise telescoping readouts now combine to give the full direct bridge identity on that same package. The inherited residual block, the completion-generated Einstein block, and the matter re-expansion block are each written as exact observer-side readouts of the ordered six-defect tuple. Summing those readouts yields \eqref{eq:direct_bridge}, and branch locking then forces every readout term to vanish on the active completed branch. Consequently
\[
\ObserverClass{\ObsBurden}=0
\]
is now proved on that branch.

The scope remains honest. This is a same-package observer-side closure theorem, not a completed first-principles substrate derivation of the Einstein sector. The next honest open burden is therefore no longer another package-preserving deeper-origin relay merely to reopen the same direct question. It is either to record an optional explicit rewritten-metric reformulation of the same closure statement, or to move beyond same-package observer-side closure and confront the genuinely deeper derivational burden of the active \Aether{} / \Aflow{} program.

\end{document}
